\section{A quantitative improvement from paired spectral powers}
\label{sec:paired-spectral-powers}

This argument improves the constant in the unrestricted cycle-count error.
It does not change its triple-logarithmic order.
The improvement uses consecutive spectral powers and a longer range of counted cycles.
The proof is mathematical and has not been formalized in Lean.

All logarithms are binary unless denoted by $\ln$.
Put $\log^{(3)}r=\log\log\log r$.

\begin{lemma}[Paired powers]\label{lem:paired-spectral-powers}
Let $J$ be a simple graph of maximum degree three, with $v$ vertices and deficit
\[
 d=\sum_{u\in V(J)}(3-\deg_J u).
\]
Let $B$ be its non-backtracking matrix.
Suppose $d<v$, and put $q=1-d/(2v)$.
If $a$ is even and $b$ is odd, with $2\le a<b$, then
\[
 \sum_{j=a}^b\frac{\operatorname{tr}B^j}{2j2^j}
 \ge q^b\sum_{j=a}^b\frac1{2j}
       -3v\sum_{j=a}^b\frac{2^{-j/2}}{2j}.
\]
\end{lemma}

\begin{proof}
The manuscript's spectral lemma gives a positive real eigenvalue at least $2-d/v$.
Every nonreal eigenvalue has modulus at most $\sqrt2$.
Every eigenvalue has modulus at most two, since the row sums of $B$ are at most two.
The trace of each power is the sum of the corresponding eigenvalue powers, with multiplicity.

Positive real eigenvalues contribute nonnegative terms.
The specified positive eigenvalue contributes at least the first term in the asserted bound.
For a negative real eigenvalue $\lambda=-2t$, where $0\le t\le1$, consecutive powers give
\[
 \frac{t^{2i}}{4i}-\frac{t^{2i+1}}{2(2i+1)}\ge0.
\]
Thus every negative real eigenvalue contributes nonnegatively over the paired interval.
There are at most $2|E(J)|\le3v$ nonreal eigenvalues.
Their total contribution is bounded below by the stated error.
\end{proof}

\begin{theorem}[An improved unrestricted forest constant]
\label{thm:paired-forest-constant}
Uniformly over connected simple graphs $G$ of rank $r$ and nonempty unions of cycles $W$ satisfying
\[
 |E(W)|\le\frac{r}{\log^2r},\qquad
 \operatorname{comp}(W)\le\frac{\log r}{100},
\]
one has
\[
 p(G,W)\le\frac12+
 \frac{1/(2\ln2)+o(1)}{\log^{(3)}r}.
\]
An explicit, deliberately generous version is
\[
 r\ge 2^{\,2^{\,2^{100}}}
 \quad\Longrightarrow\quad
 p(G,W)\le\frac12+\frac1{\log^{(3)}r}.
\]
\end{theorem}

\begin{proof}
Use the manuscript's subcubic reduction, cut lemmas, walk estimate, and second-moment variable.
Write
\[
 x=\log r,\quad y=\log x,\quad z=\log y,
 \quad D=\lfloor x/10\rfloor,\quad s=\lfloor(D-1)/2\rfloor,
\]
\[
 L=\lfloor xy/1000\rfloor,\qquad t=\lfloor y/2\rfloor.
\]
The reduction gives a subcubic graph $\Gamma$ and marked vertices $P$.
Every vertex outside $P$ has degree three, $|P|\le2r/x^2$, and
$\operatorname{comp}(\Gamma[P])\le x/100$.
Let $\mathcal F$ be the event that the uniform even subgraph is a forest on $\Gamma-P$.

First prove the explicit assertion.
Suppose $z\ge100$ and $\mathbb P(\mathcal F)>1/2+1/z$.
Choose a maximal disjoint family $\mathcal Q$ of cycles of length at most $D$ in $\Gamma-P$.
The small-cut lemma gives
\[
 |\mathcal Q|<(D+1)2^{D-1}z.
\]
Its vertex union $S$ therefore satisfies
\[
 |S|\le x^2z r^{1/10}.
\]
Put $J=\Gamma-(P\cup S)$ and $v=|V(J)|$.
Then $J$ is simple, has girth greater than $D$, and satisfies
\[
 r\le v\le3r,\qquad
 d\le\frac{6r}{x^2}+3x^2z r^{1/10}.
\]
For $z\ge100$, these estimates give
\[
 \frac{Ld}{2v}\le\frac1{100},\qquad
 \frac32 L^2 2^{-s}\le\frac1{100}.
\]
All numerical comparisons below follow from $y=2^z$, $x=2^y$, and $r=2^x$.
The threshold is intentionally much larger than necessary.

Let $a$ be the least even integer at least $6x$.
Let $b$ be the greatest odd integer at most $L$.
Lemma~\ref{lem:paired-spectral-powers} gives
\[
 \sum_{j=a}^b\frac{\operatorname{tr}B^j}{2j2^j}
 \ge\frac{99}{200}\ln(y/10000)-r^{-1}.
\]
Indeed, $a\le7x$, $b+1\ge xy/1001$, and the nonreal-eigenvalue error is at most $r^{-1}$.
Let $\mathcal C$ contain the cycles of $J$ of length at most $L$, and put
\[
 \mu=\sum_{C\in\mathcal C}2^{-|C|}.
\]
The manuscript's repeated-vertex walk estimate applies without change.
Every trace is nonnegative, so it gives
\[
 \frac{101}{100}\mu
 \ge\frac{99}{200}\ln(y/10000)-r^{-1}.
\]
For $z\ge100$, the right side implies $\mu\ge29z/100$.

For two disjoint cycles of $\mathcal C$, the manuscript's link argument gives
\[
 \ell(C,C')\le\operatorname{comp}(\Gamma[P])
  +\frac{2L}{t}+(t+1)2^{t-1}z+(8L/D)^2.
\]
For $z\ge100$, its four terms are at most $x/100,x/200,x/200,x/200$, respectively.
Consequently,
\[
 \ell(C,C')\le x/40,
 \qquad 2^{x/40}L^3 2^{-s}\le1/100.
\]
The exceptional-partner argument therefore bounds their extra contribution by $\mu/100$.
The weighted variable $Z$ from the manuscript satisfies
\[
 \mathbb EZ=\mu,\qquad
 \mathbb EZ^2\le2\mu^2+\frac{101}{100}\mu.
\]
Since $\mathcal F$ implies $Z=0$, the second-moment inequality gives
\[
 \mathbb P(\mathcal F)
 \le\frac12+\frac{101}{400\mu}
 \le\frac12+\frac{101}{116z}
 <\frac12+\frac1z.
\]
This is the required contradiction.

For the asymptotic constant, repeat the argument under
\[
 \mathbb P(\mathcal F)>\frac12+
 \frac{1/(2\ln2)+\eta}{z}
\]
for any fixed $\eta>0$.
The short-cycle deletion remains $O(x^2z r^{1/10})$.
The deficit, repeated-walk, and exceptional-partner errors now tend to zero.
The paired spectral sum gives
\[
 \mu\ge\left(\frac12-o(1)\right)\ln y.
\]
The same second moment then yields
\[
 \mathbb P(\mathcal F)
 \le\frac12+\frac{1+o(1)}{4\mu}
 \le\frac12+\frac{1/(2\ln2)+o(1)}{z},
\]
which contradicts the fixed $\eta$ margin.
The reduction transfers both conclusions to $p(G,W)$.
\end{proof}

\begin{corollary}[The unrestricted cycle constant]
\label{cor:paired-cycle-constant}
Let $\psi(r)$ be the largest number of cycles in a graph of rank $r$.
Then
\[
 \psi(r)\le2^{r-1}+
 \left(\frac1{2\ln2}+o(1)\right)
 \frac{2^r}{\log^{(3)}r}.
\]
\end{corollary}

\begin{proof}
The short-cycle edge recurrence requires the forest estimate for one cycle of length $O(\log r)$.
This length lies in the range of Theorem~\ref{thm:paired-forest-constant}.
The recurrence transfers its asymptotic constant unchanged.
\end{proof}

The original proof gives every constant larger than $2/\ln2$.
The longer cycle range gains a factor two, and the paired powers gain another factor two.
Neither change removes the logarithmic restriction imposed by the current link estimate.
